\chapter{Del modelo estático al dinámico: bancos, encaje y dolarización}
\label{cap:aplicacion}

Este capítulo hace el ejercicio completo. Toma un marco estático por capas
--- banco representativo, heterogeneidad en la exposición al choque de liquidez,
acceso desigual al interbancario, soles y dólares con respaldo público asimétrico
--- y lo traduce a un bloque de agentes heterogéneos en el sentido del capítulo
\ref{cap:bloque}, para después resolverlo en el espacio de secuencias.

El objetivo no es producir un modelo cuantitativo del sistema bancario peruano.
Es mostrar \emph{qué decisiones hay que tomar} para hacer la traducción, y qué
preguntas nuevas aparecen cuando uno la hace. La sección
\ref{sec:limitaciones-prototipo} es explícita sobre los límites del resultado.

\section{La traducción, objeto por objeto}

La tabla \ref{tab:traduccion} resume la correspondencia. Casi todo se traslada
directamente; lo que no, se señala.

\begin{table}[htbp]
\centering\small
\caption{Correspondencia entre el marco estático y el bloque dinámico.}
\label{tab:traduccion}
\begin{tabular}{@{}p{0.41\textwidth}p{0.52\textwidth}@{}}
\toprule
\textbf{Marco estático} & \textbf{Bloque en el espacio de secuencias}\\
\midrule
Tipos $k\in\{L,S\}$ con masas $\mu_k$ &
Estado exógeno $e$ con cadena de Markov $\Pi$; las masas son su distribución
estacionaria, pero ahora los bancos \emph{transitan} entre tipos\\[3pt]
Reservas $m=\rho+u$ &
Igual, pero $u^c_t$ se reelige cada período según una condición de primer orden\\[3pt]
Choque $\varepsilon_k\sim F_k$, media cero &
$\varepsilon\sim U[-a_e,a_e]$; la uniforme da forma cerrada para $G$ y $X$\\[3pt]
$\omega^c = \beta_k A^c + \varepsilon^c$, con $A^u=-\eta$ con prob. $p$ &
Estado agregado \emph{intraperiodo} $A\in\{\text{normal},\text{escaso}\}$,
integrado dentro del paso hacia atrás\\[3pt]
$\tau_k = 1-e^{-\lambda_k}$ &
$\tau(e,n_-) = \tau_e\, n_-/(n_-+\bar{n})$: acceso creciente en el tamaño\\[3pt]
$\psi_k = \tau_k\min(1,P/Q)$ &
Igual, por estado agregado: $\psi(e,n_-,A)=\tau(e,n_-)\Psi^u_t(A)$\\[3pt]
$\Delta = \sum_k\mu_k(1-\psi_k)D_k$ &
$\Delta^u_t$, agregado sobre la distribución conjunta de tipo y tamaño\\[3pt]
$F^u = v_F\rho^u D^u + R_D$, $\phi=\min(1,F^u/\Delta^u)$ &
Igual, pero $F^u_t$ y $\Delta^u_t$ son secuencias y $\phi_t$ es una incógnita del
sistema de equilibrio\\[3pt]
$C^u_k = \Ex[\phi(1-\psi_k)D^u_k]$ &
Igual; es la variable de incidencia y se calcula por tipo\\[3pt]
--- (no existe) &
Patrimonio $n_-$ como \textbf{estado endógeno}: acumulación y erosión de capital\\[3pt]
--- (no existe) &
Entrada y salida de bancos, que fija la escala del sistema\\
\bottomrule
\end{tabular}
\end{table}

Dos entradas de la tabla merecen comentario, porque no son traducciones sino
decisiones.

\paragraph{El acceso al interbancario depende del tamaño} En el marco estático
$\tau_k$ es un parámetro por tipo. Si se traslada así, la distribución de tamaños
no afecta a ningún agregado --- todos los productos del bloque serían lineales en
$n_-$ --- y el bloque heterogéneo colapsaría a tres bancos representativos, como
advertía el capítulo \ref{cap:macrofinanzas}. Hacer que $\tau$ dependa del tamaño
es lo que le da contenido a la distribución, y además conecta directamente con el
Hecho Estilizado 2 del marco estático (bancos más pequeños, depósitos más
volátiles).

\paragraph{La liquidez voluntaria se deriva, no se supone} El marco estático deja
$u^c_k$ como una decisión pendiente. Aquí se cierra con una condición de primer
orden explícita, que es lo que permite que el encaje y el colchón interactúen.

\section{El problema del banco}

\paragraph{Dentro del período} Un banco con tipo $e$ y patrimonio $n_-$ capta
depósitos $d=\lambda n_-$, de los cuales una fracción $w_e$ está en dólares.
Elige colchones voluntarios $u^s, u^u$; sus reservas son $\mathrm{res}^c=\rho^c+u^c$
y su crédito $b = n_- + d - \sum_c \mathrm{res}^c d^c$.

Después se realiza el estado agregado $A$ y el choque idiosincrático. Bajo
\emph{encaje no utilizable} (NU) la posición de liquidez por unidad de depósitos
en la moneda $c$ es $\mu^c_A = u^c + m_A$, donde $m_A$ es el desplazamiento
agregado ($0$ en el estado normal, $-\eta$ en el escaso). El faltante esperado es
\begin{equation}
G(\mu;a) = \Ex\big[(-(\varepsilon+\mu))^+\big]
= \begin{cases}
-\mu & \mu\le-a\\
(a-\mu)^2/(4a) & -a<\mu<a\\
0 & \mu\ge a
\end{cases}
\label{eq:G}
\end{equation}
y el excedente esperado es $X(\mu;a)=G(\mu;a)+\mu$, identidad que conviene usar
como prueba del código.

\paragraph{La condición de primer orden del colchón} El banco iguala el costo de
oportunidad de inmovilizar liquidez con el beneficio marginal esperado de la
cobertura:
\begin{equation}
\underbrace{r^b - r^{m,c}}_{\text{costo de oportunidad}}
=\;
\sum_A p_A\,\underbrace{\kappa_A}_{\substack{\text{costo esperado del}\\\text{faltante residual}}}
\cdot\underbrace{\Pr(\text{faltante}\mid A)}_{-G'(\mu^c_A;a)},
\qquad
\kappa_A = \iota^{\text{ef}}_A\big(1-\tau(e,n_-)\,\Psi^c_A\big),
\label{eq:cpo}
\end{equation}
con $\iota^{\text{ef}}_A = \phi_A\iota^u + (1-\phi_A)\bar\iota^u$: el faltante
cubierto por el respaldo cuesta $\iota^u$, el no cubierto cuesta $\bar\iota^u$,
que es mucho más.

\begin{idea}
La ecuación \eqref{eq:cpo} es el corazón económico del modelo y conviene leerla
despacio. El beneficio marginal de una unidad extra de colchón es la
\emph{probabilidad de que llegue a usarse}, ponderada por lo que costaría el
faltante. Esa probabilidad es $-G'$, que es continua: de ahí viene la
diferenciabilidad de todo el modelo, pese a los $\max$. Y el costo esperado
$\kappa_A$ depende de dos variables de equilibrio general --- el racionamiento
del interbancario $\Psi^u$ y la capacidad de respaldo $\phi$ --- lo que hace que
las decisiones privadas de liquidez y la capacidad pública de cobertura se
determinen \emph{simultáneamente}.
\end{idea}

Con un solo estado agregado, \eqref{eq:cpo} tiene solución cerrada,
$u^* = a(1-2\,\mathrm{spread}/\kappa)$. Con varios estados el lado derecho es
lineal a trozos y conviene resolverla por bisección: cuesta poco y evita errores
de álgebra en los casos de esquina.

\paragraph{La decisión dinámica} El resultado del período, por unidad de
patrimonio, es
\begin{equation}
R(e,n_-) = 1 + r^b + \lambda\Big[(r^b-r^d) - \kappa^{\text{op}}
- \sum_c \omega^c_e\big((r^b-r^{m,c})\,\mathrm{res}^c + \textstyle\sum_A p_A \iota^{\text{ef}}_A (1-\psi^c_A) G^c_A\big)\Big].
\end{equation}
El banquero reparte $R\,n_-$ entre dividendos y patrimonio retenido maximizando
$\Ex\sum\beta^t u(\mathrm{div}_t)$ con $u$ cóncava, y sobrevive con probabilidad
$\varsigma$; los que salen consumen su patrimonio y son reemplazados por
entrantes de tamaño fijo. Formalmente es el mismo problema de ahorro del
capítulo \ref{cap:bloque}, con un retorno que depende del tipo y del tamaño, y
se resuelve con el mismo método del punto endógeno.

\section{El grafo de equilibrio general}

Hay simultaneidad genuina: $\Psi^u$ y $\phi$ entran en \eqref{eq:cpo}, y las
decisiones que resultan de \eqref{eq:cpo} determinan la oferta, la demanda y el
faltante residual que fijan $\Psi^u$ y $\phi$. Siguiendo el capítulo
\ref{cap:dag}, el ciclo se rompe declarando ambas como incógnitas.

\begin{center}\small
\begin{tabular}{@{}lll@{}}
\toprule
\textbf{Incógnitas} & \textbf{Objetivos} & \textbf{Choques}\\
\midrule
$r^b_t$ (tasa activa) & vaciado del mercado de crédito & $\rho^u_t$ (encaje en ME)\\
$\Psi^u_t$ (holgura del interbancario) & $\Psi^u_t = \min(1, P^u_t/Q^u_t)$ & $\eta_t$ (escasez de dólares)\\
$\phi_t$ (capacidad de respaldo) & $\phi_t = \min(1, F^u_t/\Delta^u_t)$ &\\
\bottomrule
\end{tabular}
\end{center}

El grafo tiene cuatro bloques en orden topológico: \textsf{bancos} (heterogéneo),
\textsf{crédito}, \textsf{interbancario} y \textsf{respaldo} (los tres simples).
El sistema es de $3T\times 3T$ y se resuelve en segundos.

\section{El estado estacionario}

La calibración es trimestral e ilustrativa. Los encajes exigibles se fijan en los
valores agregados de julio de 2026 --- $5.56\%$ en moneda nacional y $35.37\%$ en
moneda extranjera --- y las dispersiones de los tres tipos, $a_e\in\{0.08,0.15,0.28\}$,
corresponden a desviaciones estándar del choque de retiro de $4.6\%$, $8.7\%$ y
$16.2\%$ de los depósitos, que es el rango de volatilidad del crecimiento de
depósitos documentado en el marco estático. El resto de los parámetros se eligen
para obtener un estado estacionario con las propiedades de la tabla
\ref{tab:ss}.

\begin{table}[htbp]
\centering\small
\caption{Estado estacionario del prototipo. Las tres últimas filas son los
puntos que hay que verificar antes de calcular jacobianos.}
\label{tab:ss}
\begin{tabular}{@{}lrl@{}}
\toprule
Ratio de capital & $12.50\%$ & patrimonio sobre activos\\
ROE anualizada & $12.12\%$ & dividendos sobre patrimonio\\
Colchón voluntario en ME & $16.83\%$ & de los depósitos en ME\\
Reservas totales en ME & $52.20\%$ & de los depósitos en ME\\
\midrule
$P^u/Q^u$ en el estado escaso & $0.359$ & el interbancario raciona\\
$\phi$ en el estado escaso & $0.800$ & el fondo cubre el 80\% del faltante residual\\
$v_F$ (fracción movilizable) & $0.115$ & calibrada al objetivo anterior\\
\midrule
Masa en el último punto de la grilla & $1.9\times10^{-9}$ & la grilla no ata\\
Distancia al quiebre en $\min(1,P/Q)$ & $0.64$ & el estado estacionario es interior\\
Distancia al quiebre en $\min(1,F/\Delta)$ & $0.20$ & el estado estacionario es interior\\
\bottomrule
\end{tabular}
\end{table}

\begin{figure}[htbp]
\centering
\includegraphics[width=\textwidth]{fig09_bancos_estado_estacionario.pdf}
\caption{Estado estacionario del prototipo. Izquierda: cola de la distribución de
tamaños dentro de cada tipo; la mediana está en $1.13$ y el percentil 99 en
$3.60$ veces el tamaño de entrada. Centro: el acceso al interbancario crece con
el tamaño, lo que hace que la distribución importe para los agregados. Derecha:
la incidencia estacionaria por tipo.}
\label{fig:bancos-ss}
\end{figure}

El panel derecho de la figura \ref{fig:bancos-ss} contiene el resultado
distributivo básico, que reproduce la hipótesis del marco estático y la matiza:

\begin{itemize}[leftmargin=1.4em]
\item El tipo más expuesto y con peor acceso (tipo 3) utiliza un $78\%$ más de
cobertura pública por unidad de depósito en moneda extranjera que el tipo menos
expuesto ($0.0085$ frente a $0.0048$).
\item Pero también mantiene el mayor colchón voluntario ($22.0\%$ frente a
$15.5\%$). El autoseguro \emph{atenúa} el ordenamiento sin revertirlo, que es
exactamente la salvedad que el marco estático anticipa.
\item Hay una no monotonicidad que vale la pena notar: el tipo 2 mantiene
\emph{menos} colchón que el tipo 1 ($14.2\%$ frente a $15.5\%$) pese a ser más
expuesto. La razón es que el colchón cubre dos riesgos distintos ---
el idiosincrático y el agregado --- y para un banco de baja dispersión
el motivo agregado domina y no escala con $a_e$. Es un mecanismo que el modelo
estático con un solo choque no puede generar.
\end{itemize}

\section{Los jacobianos del bloque}

\begin{figure}[htbp]
\centering
\includegraphics[width=\textwidth]{fig10_bancos_jacobianos.pdf}
\caption{Columnas del jacobiano del bloque bancario respecto del encaje exigible
en moneda extranjera. Cada curva es la respuesta a la noticia, recibida hoy, de
que $\rho^u$ será un punto más alto en la fecha $s$.}
\label{fig:bancos-J}
\end{figure}

Hay dos verificaciones analíticas que estos jacobianos pasan exactamente y que
conviene usar siempre que se construye un bloque nuevo.

\begin{enumerate}[leftmargin=1.4em]
\item $\J^{B,\rho^u}_{0,0} = -3.265699$, que coincide hasta la última cifra con
$-D^u_\sst = -3.265699$. Tiene que ser así: un punto más de encaje inmoviliza
exactamente un punto de los depósitos en moneda extranjera, y en impacto los
depósitos están predeterminados.
\item $\J^{Y,\rho^u}_{0,s} = 0$ para todo $s\ge1$ y toda variable de acervo $Y$.
Una noticia sobre el encaje futuro no puede cambiar el crédito ni el patrimonio
de hoy, porque ambos están determinados por la distribución heredada. Lo que
\emph{sí} cambia hoy es el dividendo --- el banco retiene más capital anticipando
el costo futuro --- y en efecto $\J^{\mathrm{div},\rho^u}_{0,5} = -7.8\times10^{-4}$.
\end{enumerate}

\begin{idea}
Esa pareja de verificaciones es el mejor ejemplo de lo que un jacobiano enseña.
La primera fila de la matriz dice qué puede moverse hoy: solo las variables de
flujo que son elección. Todo lo demás es acervo heredado. Si tu código produce
una respuesta no nula en la primera fila para una variable predeterminada, hay un
error de contabilidad temporal.
\end{idea}

\section{Un choque de encaje: el canal que el modelo estático no ve}

La figura \ref{fig:bancos-irf} muestra la respuesta de equilibrio general a un
aumento de un punto porcentual del encaje en moneda extranjera, con persistencia
AR(1) de $0.90$.

\begin{figure}[htbp]
\centering
\includegraphics[width=\textwidth]{fig11_bancos_irf_encaje.pdf}
\caption{Respuesta a un aumento de 1 punto porcentual del encaje exigible en
moneda extranjera ($\rho=0.90$). El patrimonio está predeterminado en impacto y
se erosiona después; el colchón voluntario \emph{cae}; la cobertura pública
utilizada \emph{sube}.}
\label{fig:bancos-irf}
\end{figure}

Los números en impacto: el crédito cae $0.32\%$, la tasa activa sube $0.31$
puntos básicos, el colchón voluntario en moneda extranjera cae $0.74\%$ y la
cobertura pública utilizada sube $3.00\%$. El patrimonio no se mueve en impacto
--- está predeterminado --- y alcanza su caída máxima de $0.10\%$ en el trimestre
12.

El resultado interesante es el del colchón voluntario, y conviene entender por
qué ocurre. Estamos en régimen de \emph{encaje no utilizable}: el requisito no
es un colchón, de modo que $\partial u^u/\partial\rho^u = 0$ \emph{dentro del
bloque} --- en efecto, $\J^{Uu,\rho^u}_{0,0}=0$ exactamente. Sin embargo, en
equilibrio general el colchón cae. El canal es indirecto:

\begin{equation}
\rho^u \uparrow
\;\Longrightarrow\;
F^u = v_F\rho^u D^u \uparrow
\;\Longrightarrow\;
\phi = \min(1, F^u/\Delta^u) \uparrow
\;\Longrightarrow\;
\iota^{\text{ef}} = \phi\iota^u + (1-\phi)\bar\iota^u \downarrow
\;\Longrightarrow\;
u^u \downarrow .
\end{equation}

Es decir: \textbf{el encaje desplaza liquidez voluntaria no porque la sustituya
mecánicamente, sino porque engorda el fondo de respaldo y abarata el faltante}.
En impacto $\phi$ sube $0.88\%$ y eso basta para que el colchón caiga $0.74\%$ y
la cobertura utilizada suba $3.00\%$. Este canal no existe en el modelo
estático con $\phi$ dado, y es el tipo de resultado que justifica el esfuerzo de
pasar a dinámico.

\paragraph{Incidencia} El panel izquierdo de la figura \ref{fig:bancos-incid}
muestra que el aumento relativo de la cobertura utilizada es mayor para el tipo
menos expuesto ($+3.97\%$) que para los tipos 2 y 3 ($+2.53\%$ y $+2.75\%$).
Conviene no leer mal este resultado: en \emph{niveles} el tipo 3 sigue usando
mucha más cobertura, porque parte de una base casi el doble de grande. Lo que
dice el ejercicio es que el ajuste marginal ante un cambio de encaje recae
proporcionalmente más sobre los bancos menos expuestos. Separar la incidencia
en niveles de la incidencia marginal es, probablemente, una distinción que vale
la pena hacer explícita en una tesis sobre incidencia.

\begin{figure}[htbp]
\centering
\includegraphics[width=\textwidth]{fig12_bancos_incidencia.pdf}
\caption{Respuesta de la cobertura pública utilizada por tipo de banco, en
porcentaje de su valor estacionario. Izquierda: choque de encaje. Derecha:
choque de escasez de dólares. El cambio de signo entre tipos en el panel derecho
es un mecanismo del modelo, no un artefacto numérico; se explica en el texto.}
\label{fig:bancos-incid}
\end{figure}

\section{Un choque de escasez de dólares}

El panel derecho de la figura \ref{fig:bancos-incid} muestra algo más llamativo.
Ante un aumento de 5 puntos porcentuales en la magnitud de la salida agregada de
dólares, la cobertura utilizada por el tipo 1 \emph{cae} un $35\%$, la del tipo 2
\emph{sube} un $12\%$ y la del tipo 3 cae un $10\%$.

El mecanismo es el siguiente, y vale la pena seguirlo porque ilustra cómo se
diagnostica un resultado contraintuitivo.

\begin{enumerate}[leftmargin=1.4em]
\item Para el tipo 1, cuya dispersión idiosincrática es pequeña ($a=0.08$), el
colchón óptimo está determinado \emph{solo} por el motivo agregado: en el estado
normal su posición de liquidez ya excede $a$ y nunca tiene faltante. La derivada
numérica confirma $\partial u^u/\partial\eta = 1.0000$ exactamente: ese banco
compensa la salida agregada uno a uno.
\item Los tipos 2 y 3 también se aseguran contra el riesgo idiosincrático, de
modo que su colchón responde solo parcialmente: $\partial u^u/\partial\eta$ vale
$0.32$ y $0.30$.
\item En equilibrio general, el choque reduce la holgura del interbancario
($-11.4\%$) y la capacidad de respaldo ($-6.8\%$), lo que encarece el faltante
residual y empuja a todos los bancos a aumentar su colchón más allá de la
compensación mecánica. El colchón agregado sube $25.8\%$.
\item Para el tipo 1, ese exceso de autoseguro domina: su posición de liquidez
esperada \emph{mejora}, su faltante esperado cae y termina usando menos
cobertura pública. Para el tipo 2 no alcanza y termina usando más.
\end{enumerate}

\begin{idea}
La lectura económica es que un choque de escasez de dólares provoca una
\emph{huida hacia el autoseguro privado}, y que esa huida es más intensa
--- en términos relativos --- entre los bancos cuyo único motivo de liquidez es
el riesgo agregado. El resultado es una recomposición de quién usa el respaldo
público, no solo un cambio de nivel. Es exactamente la clase de pregunta
distributiva que motiva el marco estático, y solo el modelo dinámico con
equilibrio general simultáneo la puede contestar.
\end{idea}

\section{Encaje utilizable y el umbral}
\label{sec:regimenes}

El marco estático menciona, como resultado secundario, que bajo \emph{encaje
utilizable} un aumento del requisito podría sustituir liquidez voluntaria. El
prototipo lo hace preciso.

\begin{figure}[htbp]
\centering
\includegraphics[width=\textwidth]{fig13_bancos_regimenes.pdf}
\caption{Colchón voluntario óptimo en función del encaje exigible, por tipo de
banco. Bajo encaje no utilizable (izquierda) el requisito no afecta la decisión
privada. Bajo encaje utilizable (derecha) la sustituye uno a uno hasta un umbral
--- el nivel que el banco elegiría por sí mismo --- y después deja de tener
efecto marginal.}
\label{fig:regimenes}
\end{figure}

Bajo encaje utilizable, el colchón efectivo es $\mu^c = \max\{\rho^c, \mu^{c*}\}$,
donde $\mu^{c*}$ resuelve \eqref{eq:cpo}. Por lo tanto
\begin{equation}
u^c = \max\{\mu^{c*}-\rho^c,\,0\},
\qquad
\frac{\partial u^c}{\partial \rho^c} =
\begin{cases}
-1 & \text{si } \rho^c < \mu^{c*} \quad\text{(sustitución uno a uno)}\\
0 & \text{si } \rho^c > \mu^{c*} \quad\text{(el requisito es holgado)}
\end{cases}
\end{equation}

El umbral $\mu^{c*}$ depende del tipo: en la calibración vale $0.159$, $0.144$ y
$0.223$ para los tres tipos. Con un encaje exigible en moneda extranjera de
$35.4\%$, los tres están muy por encima de su umbral, de modo que bajo encaje
utilizable el requisito eliminaría prácticamente todo el riesgo de liquidez en
dólares --- y con él, el motivo asegurador que el marco estático quiere estudiar.

\begin{tutesis}
Esto es un resultado con consecuencias directas para tu diseño de investigación,
y conviene tenerlo claro antes de IE2. La distinción entre encaje utilizable y
no utilizable no es un detalle de implementación: determina si el instrumento
\emph{tiene} un margen interesante. Con $\rho^u=35\%$ y choques de retiro de
hasta $28\%$ de los depósitos, bajo encaje utilizable no hay faltante en casi
ningún estado y el modelo se vuelve degenerado. Tu interpretación del encaje
como un \emph{fondo} que financia un seguro --- es decir, el régimen NU --- no es
solo la más fiel al arreglo institucional peruano: es la única bajo la cual la
pregunta de incidencia tiene contenido con encajes de esa magnitud. Vale la pena
decirlo así, con el número al lado, en lugar de presentar NU como un supuesto
auxiliar.
\end{tutesis}

\section{Qué habría que arreglar antes de usar esto cuantitativamente}
\label{sec:limitaciones-prototipo}

El prototipo es correcto como ilustración del método y es honesto decir dónde no
alcanza.

\begin{enumerate}[leftmargin=1.4em]
\item \textbf{La fricción de pago de dividendos es una forma reducida.} La
utilidad cóncava sobre dividendos es un sustituto conveniente de un costo de
emisión de capital o de un valor de franquicia, pero no está microfundada y sus
parámetros no están disciplinados por nada. Para un ejercicio cuantitativo
convendría reemplazarla por una especificación tipo Gertler--Karadi con una
restricción de incentivos explícita.
\item \textbf{La restricción de apalancamiento es exógena y siempre ata.} En los
datos el ratio de capital varía por banco y en el tiempo, y a veces no ata.
\item \textbf{Los choques en soles y en dólares son independientes.} En una
crisis de dolarización están correlacionados, y esa correlación es probablemente
importante para la demanda de cobertura conjunta.
\item \textbf{La composición monetaria $w_e$ es fija.} El marco estático lo
justifica como aproximación de primer orden, y es razonable; pero en un modelo
dinámico la sustitución de depósitos entre monedas ante un cambio de encaje es un
canal de política de primer orden.
\item \textbf{No hay precio del dólar.} Un modelo que quiera hablar de escasez de
dólares en serio necesita un tipo de cambio y una condición de arbitraje.
\item \textbf{La calibración es ilustrativa.} Las dispersiones $a_e$ están
elegidas para abarcar el rango observado de volatilidad de depósitos, no
estimadas; $\tau_e$, $\iota^u$ y $\bar\iota^u$ no están disciplinados por nada.
El capítulo \ref{cap:hoja-ruta} propone cómo disciplinarlos.
\end{enumerate}

Ninguna de estas limitaciones es del método: todas son del modelo. Esa distinción
importa, porque significa que arreglarlas no obliga a cambiar de herramienta.
